\section{A new function on square matrices}

In the previous chapter we learned how to read matrices and how to perform basic operations on them. We now introduce a new object associated with a square matrix: a single number called its \emph{determinant}.

For each size \(n\), the determinant is a function
\[
\det\colon \mathbb{R}^{n\times n}\longrightarrow \mathbb{R}.
\]
It takes an \(n\times n\) matrix as input and returns one real number.

At this point we should not pretend that the formula is forced by ideas that have not yet been developed. We will begin by defining determinants explicitly for matrices of size \(1\times1\), \(2\times2\), and \(3\times3\). The definition may initially appear somewhat artificial. Its usefulness will become visible through the properties that follow from it.

This is a common order in mathematics:
\[
\text{definition}\quad\longrightarrow\quad
\text{calculation}\quad\longrightarrow\quad
\text{properties}\quad\longrightarrow\quad
\text{interpretation}.
\]
We first agree what the object is. We then investigate what the definition allows us to do.

\begin{remarkbox}
The determinant is defined only for square matrices in this course. Thus expressions such as
\[
\det\begin{pmatrix}1&2&3\\4&5&6\end{pmatrix}
\]
are not defined.
\end{remarkbox}

The chapter will follow a controlled progression.
\begin{itemize}
\item We define determinants of small matrices and practise direct calculation.
\item We examine how the value changes when rows or columns are modified.
\item We study what can be concluded when the determinant is zero.
\item We extend the definition to arbitrary square matrices using minors, cofactors, and Laplace expansion.
\item We collect general properties that make determinant calculations efficient.
\end{itemize}

The main question is therefore not yet ``What deep object does the determinant represent?'' The immediate question is more elementary and more honest:

\begin{summarybox}
How is the determinant defined, and what follows from that definition?
\end{summarybox}
